\section*{Capítulo \ref{ch:b2:viscous-flows} --- Flujos viscosos}

\begin{solution}{exo:b2:viscous-flows:1}
$\tau = \eta U/h = 0.25 \times 2/2 \times 10^{-4} = \qty{2.5}{kPa}$; $F = \qty{1.25}{kN}$;
$P = Fv = \qty{2.5}{kW}$. Rozamiento seco: $0.3 \times 98 = \qty{29}{N}$ --- aquí la
película fina y muy cizallada resiste mucho más; el arrastre viscoso crece
como $U/h$, de modo que la lubricación compensa a velocidades menores y
con películas más gruesas (y nunca se agarrota).
\end{solution}

\begin{solution}{exo:b2:viscous-flows:2}
(a) $2 \times 10^{-6} \times 3 \times 10^{-5}/10^{-6} = 6 \times 10^{-5}$, reptante. (b)
$1.8 \times 1.5/10^{-6} = 2.7 \times 10^6$, turbulento. (c) $4 \times 30/1.5 \times 10^{-5}
= 8 \times 10^6$, turbulento. (d) $0.025 \times 0.33/3.3 \times 10^{-6} = 2500$,
laminar al límite. (e) $900 \times 0.5 \times 0.01/0.1 = 45$, laminar.
\end{solution}

\begin{solution}{exo:b2:viscous-flows:3}
$D_V = \pi R^4\Delta P/8\eta L = \pi(2 \times 10^{-4})^4 \times 2 \times 10^4/(8 \times 10^{-3}
\times 0.03) = \qty{0.42}{mL/s}$; \qty{12}{s} para \qty{5}{mL}. $v_{\max} = \Delta PR^2/
4\eta L = \qty{6.7}{m/s}$, $\langle v\rangle = \qty{3.3}{m/s}$, $\mathrm{Re} = 2R\langle v
\rangle/\nu = 1300$: laminar. $\tau_w = \Delta PR/2L = \qty{67}{Pa}$. Con \qty{0.8}{mm}:
dieciséis veces el caudal, \qty{6.7}{mL/s} --- pero $\mathrm{Re} \approx 10^4$: el flujo
se vuelve turbulento y la ganancia es menor.
\end{solution}

\begin{solution}{exo:b2:viscous-flows:4}
$v_\infty = 2R^2\Delta\rho\,g/9\eta$. (a) \qty{1.3}{cm/s}, $\mathrm{Re} = 0.02$, \qty{75}{s}
por metro. (b) \qty{0.24}{mm/s}, $\mathrm{Re} = 3 \times 10^{-5}$, \qty{70}{min} por
metro. (c) \qty{1.2}{m/s} pero $\mathrm{Re} = 16$: Stokes sobreestima (la
celeridad real es de unos \qty{0.7}{m/s}), cosa de un segundo por metro.
\end{solution}

\begin{solution}{exo:b2:viscous-flows:5}
(a) $\langle v\rangle = 0.1/0.0707 = \qty{1.4}{m/s}$, $\mathrm{Re} = 4.2 \times 10^5$. (b)
$8\eta LD_V/\pi R^4 = \qty{500}{Pa}$. (c) $\lambda = 0.3/25.5 = 0.0118$, $\Delta P = 0.0118
\times 3333 \times 1000 = \qty{39}{kPa}$: $80$ veces más; $P = \Delta PD_V = \qty{3.9}{kW}$.
(d) Misma celeridad, $\mathrm{Re} = 3 \times 10^5$, $\lambda = 0.0128$, $\Delta P = 0.0128 \times 4717
\times 1000 = \qty{60}{kPa}$, \qty{6}{kW}: más pared por unidad de caudal, más pérdida.
\end{solution}

\begin{solution}{exo:b2:viscous-flows:6}
(a) $v = 2 \times 25 \times 10^{-12} \times 50 \times 9.81/(9 \times 10^{-3}) = \qty{2.7}{\micro m/s}$;
\qty{5}{cm} en $1.8 \times 10^4$ s, cinco horas. (b) \qty{2.7}{mm/s}, \qty{18}{s}. (c)
$\Omega = \sqrt{1000g/r} = \qty{313}{rad/s} = \qty{3000}{rpm}$. (d) $v \propto R^2$: las
partículas grandes sedimentan primero; giros sucesivos a velocidad
creciente separan los tamaños.
\end{solution}

\begin{solution}{exo:b2:viscous-flows:7}
(a) $F = 0.6 \times 144 \times 0.35 = \qty{30}{N}$, $+\qty{4}{N}$: $P = 34 \times 12 =
\qty{410}{W}$. (b) $47 + 4 = \qty{51}{N}$, \qty{770}{W}. (c) Celeridad relativa
\qty{17}{m/s}: $61 + 4 = \qty{65}{N}$, por la velocidad respecto del suelo: \qty{780}{W}. (d)
Resistencia $\qty{21}{N}$: \qty{300}{W}, con un ahorro de \qty{110}{W}. (e) $mg\sin\alpha = \qty{39}{N} =
4 + 0.21v^2$: $v = \qty{13}{m/s}$, \qty{47}{km/h}.
\end{solution}

\begin{solution}{exo:b2:viscous-flows:8}
(a) Golf: $\tfrac12\rho C_xS = \qty{2.2e-4}{kg/m}$, $v_\infty = \sqrt{0.45/2.2 \times 10^{-4}}
= \qty{45}{m/s}$; paracaidista $\sqrt{785/0.42} = \qty{43}{m/s}$; con el paracaídas abierto $\sqrt{785/15} =
\qty{7.2}{m/s}$. (b) $m\dot v = mg - kv^2$, es decir $\dot v = g(1 - v^2/v_\infty^2)$,
cuya solución desde el reposo es $v_\infty\tanh(gt/v_\infty)$. (c) $\tanh x = 0.9$ en
$x = 1.47$: $t = 1.47 \times 43/9.81 = \qty{6.5}{s}$; $z = (v_\infty^2/g)\ln\cosh x =
190 \times 0.83 = \qty{160}{m}$.
\end{solution}

\begin{solution}{exo:b2:viscous-flows:9}
(a) Velocidad de cizalla $\Omega R/e = 0.5/10^{-3} = \qty{500}{s^{-1}}$; $\Gamma = \tau\cdot2\pi Rh
\cdot R$: $\tau = 0.02/(2\pi \times 2.5 \times 10^{-3} \times 0.1) = \qty{12.7}{Pa}$. (b) $\eta
= 12.7/500 = \qty{2.5e-2}{Pa.s}$. (c) $\mathrm{Re} = \rho Ue/\eta \approx 1000 \times 0.5 \times
10^{-3}/0.025 = 20$: laminar. (d) Un hueco estrecho hace uniforme la
velocidad de cizalla (la aproximación plana) y mantiene laminar el flujo.
\end{solution}

\begin{solution}{exo:b2:viscous-flows:10}
$D_V = \qty{8.3e-5}{m^3/s}$. (a) $R_h = 8\eta L/\pi R^4 = 8 \times 3.5 \times 10^{-3} \times
0.4/(\pi \times 2.4 \times 10^{-8}) = \qty{1.5e5}{Pa.s/m^3}$, $\Delta P = \qty{12}{Pa}$. (b)
Un capilar, $\qty{3.5e16}{Pa.s/m^3}$; $10^{10}$ en paralelo: \qty{3.5e6}{Pa.s/m^3};
$\Delta P = \qty{290}{Pa}$ (\qty{2}{mmHg}). (c) Una arteriola, $8 \times 3.5 \times 10^{-3}
\times 2 \times 10^{-3}/(\pi \times 10^{-20}) = \qty{1.8e15}{Pa.s/m^3}$; el conjunto,
\qty{1.8e7}{Pa.s/m^3}; $\Delta P = \qty{1.5}{kPa}$: las arteriolas dominan sobre
las otras dos (la presión cae en los vasos pequeños, no en los grandes).
(d) $P = 13000 \times 8.3 \times 10^{-5} = \qty{1.1}{W}$, alrededor del $1\%$. (e)
$(1/1.1)^4 = 0.68$: se elimina un tercio de la resistencia arteriolar ---
así bajan la presión arterial los vasodilatadores.
\end{solution}

\begin{solution}{exo:b2:viscous-flows:11}
(a) $F = b\int_0^L\tau_w\,\dd x = 0.332\rho U^2b\sqrt{\nu/U}\int_0^L\dd x/\sqrt x =
0.664\rho U^2b\sqrt{\nu L/U} = 0.664\rho U^2bL/\sqrt{\mathrm{Re}_L}$. (b) $\mathrm{Re}_L =
3 \times 10^5$ (laminar); $\delta \approx 5\sqrt{\nu L/U} = \qty{9}{mm}$; $F = 0.664 \times 1000
\times 0.09 \times 1/548 = \qty{0.11}{N}$ por cara, \qty{0.22}{N}. (c) $\mathrm{Re}_L = 10^7$:
la capa es turbulenta más allá de los primeros centímetros y la fórmula
laminar falla (la resistencia real es varias veces mayor). (d) $\mathrm{Re}_x =
10^9$: $\delta \approx 0.37 \times 100/63 = \qty{0.6}{m}$.
\end{solution}

\begin{solution}{exo:b2:viscous-flows:12}
(a) $0 = \rho g + \eta v''(y)$, $v(0) = 0$ (adherencia), $v'(h) = 0$ (sin esfuerzo en
la superficie libre). (b) Integrando dos veces: $v = (\rho g/2\eta)(2hy - y^2)$; $q =
\int_0^hv\,\dd y = \rho gh^3/3\eta$. (c) $v_s = \rho gh^2/2\eta = 1200 \times 9.81 \times
4 \times 10^{-8}/1 = \qty{0.47}{mm/s}$; $q = \qty{6.3e-8}{m^2/s}$; \qty{28}{cm} en diez
minutos --- una pintura newtoniana de ese espesor se escurriría; las
pinturas reales son seudoplásticas y tixotrópicas, y aun así se descuelgan
cuando se aplican demasiado espesas, porque $v_s \propto h^2$. (d) $1400 \times 9.81 \times 4 \times 10^{-6}/20 = \qty{2.7}{mm/s}$.
\end{solution}

\begin{solution}{pb:b2:viscous-flows:1}
\textbf{1.} $S = \qty{4.9}{cm^2}$, $\langle v\rangle = 83/4.9 = \qty{0.17}{m/s}$; $\mathrm{Re}
= 1060 \times 0.17 \times 0.025/3.5 \times 10^{-3} = 1300$: laminar (los picos
pulsátiles llegan a $\sim 5000$).

\textbf{2.} $v_{\max} = \qty{0.34}{m/s}$; $\tau_w = 4\eta\langle v\rangle/R = \qty{0.19}{Pa}$.

\textbf{3.} $\mathrm{Re} = 1060 \times 3 \times 10^{-4} \times 8 \times 10^{-6}/3.5 \times 10^{-3} =
7 \times 10^{-4}$; $\Delta P = 8\eta L\langle v\rangle/R^2 = 8 \times 3.5 \times 10^{-3} \times 10^{-3}
\times 3 \times 10^{-4}/1.6 \times 10^{-11} = \qty{525}{Pa} = \qty{3.9}{mmHg}$.

\textbf{4.} $1/0.3 = \qty{3.3}{s}$: lo bastante para que el oxígeno difunda
los pocos micrómetros que lo separan del tejido
(\cref{ch:b2:particle-diffusion}).

\textbf{5.} Una arteriola, $8\eta L/\pi R^4 = \qty{1.8e15}{Pa.s/m^3}$; $10^8$ en
paralelo: \qty{1.8e7}{Pa.s/m^3}; $\Delta P = 1.8 \times 10^7 \times 8.3 \times 10^{-5} =
\qty{1.5}{kPa} = \qty{11}{mmHg}$.

\textbf{6.} $P = \Delta P\,D_V = 13000 \times 8.3 \times 10^{-5} = \qty{1.1}{W}$, alrededor
del $1\%$ de la potencia del cuerpo.

\textbf{7.} $(1/0.6)^4 = 7.7$. En reposo los vasos de aguas abajo se dilatan
y compensan; durante el esfuerzo la demanda se multiplica por cuatro o
cinco y la arteria estrechada no puede satisfacerla: angina de pecho.

\textbf{8.} Cinco veces la celeridad: $\mathrm{Re} \approx 6500$, turbulento --- los
soplos que oye el estetoscopio durante el ejercicio.

\textbf{9.} En el flujo lento y poco cizallado de las venas pequeñas y los
capilares los glóbulos rojos se agregan y la \omterm{def:b2:viscous-flows:viscosity}{viscosidad} aparente sube (lo
contrario ocurre en la aorta, rápida); todas las estimaciones de caída de
presión en los vasos pequeños son cotas inferiores.

\textbf{10.} $\langle v\rangle = 3/1.13 = \qty{2.65}{m/s}$; $\mathrm{Re} = 850 \times 2.65 \times
1.2/0.02 = 1.35 \times 10^5$: turbulento.

\textbf{11.} $\lambda = 0.316/19.2 = 0.0165$; $\Delta P/L = (0.0165/1.2) \times \tfrac12 \times
850 \times 7.0 = \qty{41}{Pa/m} = \qty{41}{kPa/km}$; \qty{534}{bar} en toda la línea.

\textbf{12.} $534/80 = 6.7$: al menos siete estaciones.

\textbf{13.} $P = \Delta P\,D_V = 5.34 \times 10^7 \times 3 = \qty{160}{MW}$;
\qty{53}{MJ/m^3}, el $0.15\%$ de la energía del crudo.

\textbf{14.} $\Delta P = 32\eta L\langle v\rangle/D^2 = 32 \times 0.02 \times 1.3 \times 10^6
\times 2.65/1.44 = \qty{15}{bar}$: treinta y cinco veces menos. La turbulencia
es cara; unas pocas partes por millón de polímeros largos la retrasan y
recortan la caída en decenas de por ciento.

\textbf{15.} $\mathrm{Re} = 1350$: laminar; Poiseuille: $\Delta P = 32 \times 2 \times 1.3 \times
10^6 \times 2.65/1.44 = \qty{1500}{bar}$ --- imposible; el crudo debe mantenerse
caliente, de ahí el aislamiento y el calentamiento.

\textbf{16.} $1.3 \times 10^6/2.65 = \qty{4.9e5}{s}$, es decir $5.7$ días.

\textbf{17.} $\tau_w = (\Delta P/L)D/4 = 41 \times 0.3 = \qty{12}{Pa}$; $F = \tau_w\pi DL =
\qty{6e7}{N}$.

\textbf{18.} $\Delta P \propto \lambda\langle v\rangle^2/D \propto D^{-4.75}$: $1.2^{4.75} = 2.4$.

\textbf{19.} $v = 2R^2\Delta\rho\,g/9\eta = 2 \times 2.5 \times 10^{-9} \times 1800 \times 9.81/
0.18 = \qty{0.49}{mm/s}$; $\mathrm{Re} = 2 \times 10^{-3}$; \qty{41}{min} para
atravesar \qty{1.2}{m}. Las fluctuaciones turbulentas, de unos centímetros
por segundo, lo mantienen en suspensión mientras el crudo circula; en una
línea parada sedimenta en menos de una hora.

\textbf{20.} $v = 2 \times 2.5 \times 10^{-11} \times 1000 \times 9.81/1.62 \times 10^{-4} =
\qty{3.0}{mm/s}$; $\mathrm{Re} = 2 \times 10^{-3}$; \qty{1}{km} en $3.3 \times 10^5$ s, cuatro
días: cualquier corriente ascendente de más de unos milímetros por segundo la sostiene.

\textbf{21.} \qty{1.2}{m/s}, $\mathrm{Re} = 16$: Stokes sobreestima en un factor
dos aproximadamente (real: \qty{0.7}{m/s}).

\textbf{22.} Stokes daría \qty{120}{m/s} y $\mathrm{Re} \sim 10^4$: absurdo.
Resistencia cuadrática, $\tfrac43\pi R^3\rho_wg = \tfrac12\rho v^2\pi R^2C_x$: $v = \sqrt{8R
\rho_wg/3\rho C_x} = \qty{6.6}{m/s}$ ($\mathrm{Re} = 900$, coherente); \qty{150}{s} desde
\qty{1}{km}.

\textbf{23.} \qty{1}{cm/s} $>$ \qty{3}{mm/s}: la gotita sube; la gota de
llovizna necesita alrededor de \qty{1}{m/s}, que las nubes convectivas proporcionan.

\textbf{24.} $2\gamma/R$: \qty{29}{kPa}, \qty{1.4}{kPa}, \qty{144}{Pa}; presiones
dinámicas $\tfrac12\rho v^2$: $5 \times 10^{-6}$, $0.3$ y \qty{26}{Pa}. El cociente
llega a $0.2$ a \qty{2}{mm}; a \qty{5}{mm} ($v \approx \qty{9}{m/s}$, $\tfrac12\rho v^2 =
\qty{49}{Pa}$ frente a \qty{58}{Pa}) la presión aerodinámica achata la gota
y la rompe: ninguna gota de lluvia pasa de unos \qty{5}{mm}.

\textbf{25.} Aorta $10^3$ (laminar, cerca de la transición); capilar $10^{-3}$
(Poiseuille); arteriola $10^{-2}$ (Poiseuille); oleoducto caliente $10^5$
(turbulento, $\lambda(\mathrm{Re})$); oleoducto frío $10^3$ (Poiseuille); arena en
crudo $10^{-3}$ (Stokes); gotita de nube $10^{-3}$ (Stokes); gota de lluvia $10^3$
(resistencia cuadrática, $C_x \approx 0.5$).
\end{solution}
